% ch5_impedance_locus.tex --- Apparent impedance locus on R-X plane
\begin{tikzpicture}
\begin{axis}[
  thesis style,
  width  = 0.68\textwidth,
  height = 7.0cm,
  xlabel = {Resistance $R$ ($\Omega$)},
  ylabel = {Reactance $X$ ($\Omega$)},
  xmin=-5, xmax=35, ymin=-5, ymax=30,
  xtick={0,5,10,15,20,25,30},
  ytick={0,5,10,15,20,25},
  legend pos=north west,
]

%% True fault locus (no IBR): linear from origin at line impedance angle
\addplot[thick, thesisgreen, domain=0:10, samples=10]
  {1.2*x};
\addlegendentry{True fault locus ($\kibr = 0$)};

%% IBR locus: kibr = 0.3, lagging reactive injection (overreach)
\addplot[thick, thesisred, domain=0:10, samples=10]
  {1.2*x + 3};
\addlegendentry{Apparent impedance ($\kibr = 0.3$, lagging)};

%% IBR locus: kibr = 0.55, more overreach
\addplot[thick, thesisorange, domain=0:10, samples=10]
  {1.2*x + 5.5};
\addlegendentry{Apparent impedance ($\kibr = 0.55$, lagging)};

%% Zone 1 boundary (simplified rectangle)
\addplot[dashed, thick, thesisblue, domain=0:9, samples=2]
  {11.5};
\node[font=\tiny, thesisblue] at (axis cs:4, 12.5) {Zone~1 boundary};

\end{axis}
\end{tikzpicture}
